\newpage

\subsubsection{隐零点}

\begin{example}
    （25 高三上•上海杨浦•期末）已知函数 $\boldsymbol{u}(\boldsymbol{x})=\boldsymbol{e}^{x+m}, \boldsymbol{v}(\boldsymbol{x})=\boldsymbol{l} \boldsymbol{n} \boldsymbol{x}-\boldsymbol{m}(\boldsymbol{m} \in \boldsymbol{R})$ ，设 $f(x)=u(x)+v(x)$ ， $g(x)=u(x)-v(x)$.
    \begin{itemize}
        \item 当 $m=0$ 时，求曲线 $y=f(x)$ 在点 $(1, \mathrm{e})$ 处的切线方程；
        \item 证明：当函数 $y=f(x)$ 经过点 $(1,2)$ 时函数 $y=g(x)$ 有且仅有一个零点；
        \item 证明：对 $m<-1$ ，存在实数 $n$ 使 $|\boldsymbol{g}(\boldsymbol{x})|=\boldsymbol{n}$ 恰有三个不同的实数根，并指出实数 $n$ 的取值范围。
    \end{itemize}
\end{example}

\begin{tips}
    这里需要找出 $x_0 \in (1, -m)$ 的范围内.
\end{tips}

\v{20em}

